\subsection{Proof of closedness of the general integral graph}
\label{sec:closure-proof}
\begin{proof}[Proof of Theorem~\ref{thm:graph-closure}]
\emph{Choose one original-price approximation.}
The common null event on which the price bound fails belongs to the initial sigma-field by the usual conditions. Set the price to zero there; this gives an everywhere bound without changing gains up to indistinguishability. Apply bounded truncation and elementary density from
Theorem~\ref{thm:integral-input} to each $Y_n$.
Diagonalize with geometrically small test-uniform errors on integer horizons
to obtain elementary original-price gains $P_n$ converging to the same $X$
in Emery topology. By assumption $X$ is zero-start, strongly adapted and
c\`adl\`ag. Proposition~\ref{prop:gi-limit} makes it a good integrator, and
Theorem~\ref{thm:regular-gi-decomposition} gives a decomposition in the domain
of the auxiliary topology. No integral representation of $X$ is assumed.

\emph{Localize the same sequence.}
Apply Proposition~\ref{prop:topology-localization} to $P_n-X$ to obtain
vanishing auxiliary $r^1$ cost. Its topology comparison uses completeness
of the space of decomposable processes and a closed-graph theorem,
independently of integral-range closedness. Take a faster subsequence with
summable adjacent costs. Proposition~\ref{prop:joint-costs} chooses this
subsequence, one equivalent measure $Q$, and one stopping sequence $\tau_r$
so that prelocal component costs and expected boundary jumps are jointly
summable. It retains $\tau_r\leq r+1$ and $\tau_r\to\infty$.

\emph{Pass the bounds to actual stopped integral components.}
Each $P_n^{\tau_r}$ is an already realized original-price integral.
Lemma~\ref{lem:canonical-cost} and equation~\eqref{eq:actual-closed-cost}
transfer the bounds obtained from arbitrary decompositions to the canonical
components of this same gain. Subsection~\ref{sec:actual-series} applies the
bounds to every adjacent difference at the same $r$ and constructs component
series. Their limits are locally uniform on the M side and in variation on
the FV side; their sum is $X^{\tau_r}$. No unrelated sequence or independently
selected decomposition is substituted.

\emph{Realize both components with one coefficient.}
The joint stops of Subsection~\ref{sec:joint-realization} give terminal
$L^2(Q)$ bounds on the approximating M components, uniform in their index.
Together with summable adjacent FV variations, these are precisely the
hypotheses of Lemma~\ref{lem:joint-coefficient}. That lemma matches the
coefficient limit under the FV variation control with the energy limit
of convex averages of the same coefficients. One predictable coefficient
therefore represents both stopped components, whose sum is the same
$X^{\tau_r}$. The FV control also determines the coefficient where the energy
measure vanishes; no absolute continuity between the two measures is needed.

\emph{Return to the original measure and all time.}
Truncation inclusion and measure transfer from
Theorem~\ref{thm:integral-input} return each stopped general graph to the
original measure. Apply the stopping-interval pasting proved in
Section~\ref{sec:integral-calculus} to this stopping sequence.
Joining coefficients on disjoint intervals, equality of the stopped gains
and $\tau_r\to\infty$ give the all-time original-price graph $(H,X)$.
Pasting needs indistinguishable gains, not agreement of independently chosen
coefficients on overlaps. The auxiliary special decomposition is not itself
transferred between measures.
\end{proof}

\paragraph{Gain domains and terminal sets}
An elementary-test realization consists of elementary predictable strategies
with a deterministic bound on the sum of coefficient magnitudes in each row,
and a strongly adapted c\`adl\`ag ucp limit of their gains, where the gains
are Cauchy uniformly over elementary tests.
\begin{proposition}[Identification of gains and terminal claims]
For the same bounded source, the gain domains of the general truncation graph
and elementary-test realizations agree up to indistinguishability.
Their $a$-admissible terminal claim sets, defined using each candidate's own
finite terminal limit, agree as well.
\end{proposition}
\begin{proof}
Each elementary row has bounded coefficient, hence belongs to the general
graph. Its Émery Cauchy property and ucp limit give Émery convergence to that
same limit; apply Theorem~\ref{thm:graph-closure}.
Conversely approximate each bounded truncation row by elementary integrals
of the original source. Removing a stop costs at most its probability of
occurring before the horizon, independently of the test. Combine this with
truncation convergence. At integer horizon $n+1$ choose error at most
$1/(n+1)$. Horizon monotonicity gives convergence on every finite horizon,
then the common test Cauchy bound and zero start give the required ucp limit.
Indistinguishability transfers all-time lower bounds and the candidate's own
terminal limit. Existence of a terminal limit is never inferred from
finite-horizon convergence alone.
\end{proof}
